\documentclass[conference]{IEEEtran}
\usepackage[utf8]{inputenc}
\usepackage[T1]{fontenc}
\usepackage{cite}
\usepackage{amsmath,amssymb}
\usepackage{graphicx}
\usepackage{booktabs}
\usepackage{url}

\title{Effectiveness of a Generate--Verify--Repair Pipeline with Batfish for LLM‑Generated Vendor CLI Configurations}

\author{
\IEEEauthorblockN{Autonomous AI Researcher}
\IEEEauthorblockA{CNSM 2026 Agentic AI Researcher Track}
}

\begin{document}

\maketitle

\begin{abstract}
We report a preregistered paired experiment (cnsms2026-001-gvr-batfish-prereg-v1) that evaluated whether a generate--verify--repair (GVR) pipeline---shared initial LLM candidate (gpt-5-mini) + a deterministic Batfish‑style validator + at most one automated LLM repair---reduces deterministic‑verifier detected semantic/safety failures on intent\ensuremath{\rightarrow}vendor‑CLI tasks. Planned N=300 pairs; 266 complete pairs were analysed after preregistered exclusions. Baseline pass-rate = 260/266 (97.744\%); guarded (final GVR) pass-rate = 265/266 (99.624\%). Paired difference = 0.01879699248; exact McNemar two‑sided p = 0.0625. Paired bootstrap (B=10,000, seed=1729) 95\% percentile CI = [0.0037593984962406013, 0.03759398496240601]. Repair was invoked for 6 tasks; median repair iterations per task = 0. Per‑episode latency was not recorded in per‑task artifacts and thus could not be evaluated. Under shared\_initial\_candidate semantics the observed absolute improvement is small and did not meet the preregistered \ensuremath{\alpha}=0.05 decision rule. We provide preregistered design details, deterministic validator semantics, exact paired counts, representative artifact‑grounded repair examples, contamination findings (no exact public duplicates flagged), and a focused discussion of limitations and implications.
\end{abstract}

\section{Introduction}
Generative large language models (LLMs) are actively explored for NetOps tasks such as intent\ensuremath{\rightarrow}CLI translation, zero‑touch configuration synthesis, and automated maintenance workflows [1], and surveys emphasize both opportunity and risk when LLMs are applied to operational network configuration [2]. Stochastic LLM outputs can produce syntactic or semantic violations---missing required commands, unsafe command orderings, or constraint breaches---that create operational risk if applied directly. To mitigate that risk, pipelines that interpose deterministic verifiers (e.g., Batfish or header‑space analyses) and bounded repair loops---collectively generate--verify--repair (GVR)---have been proposed to provide deterministic acceptance criteria and automated correction [3,4]. Prior work presents components and prototypes, but preregistered, paired evaluations that archive verifier traces and quantify verifier‑triggered repair benefit on reproducible NetOps task banks remain limited.

This study tested the preregistered question (cnsms2026-001-gvr-batfish-prereg-v1): Does a GVR pipeline that uses a single sampled initial LLM candidate per task (shared\_initial\_candidate semantics), a deterministic Batfish‑style validator, and at most one automated LLM repair call reduce deterministic‑verifier detected semantic/safety failures relative to the same initial candidate without repair? Our principal contribution is a preregistered, artifact‑archived paired evaluation executed under the CNSM Agentic AI Researcher constraints that reports exact paired counts, inferential outcomes, execution accounting, representative artifact examples, and an evidence‑grounded interpretation.

\section{Related Work}
LLM applications in NetOps include intent\ensuremath{\rightarrow}CLI translation systems and focused benchmarks studying automated configuration synthesis and correctness [1]. Surveys synthesize common failure modes (including hallucination and constraint violations) and advocate grounding and verification components for automation [2]. Generate--verify--repair architectures have been proposed and explored in code generation and regulated domains to achieve deterministic validation via external oracles and bounded repair loops [3]. Network verification tools such as Batfish and header‑space analyses provide deterministic reachability and ordering checks and are natural verifier choices in LLM‑augmented NetOps prototypes; prior analyses document Batfish's design and motivate its integration in automated pipelines [4]. Work examining what LLMs need to synthesize correct router configurations highlights sequencing and vendor‑syntax difficulties that verifiers can detect and flag for repair [5].

Our study differs by executing a preregistered paired test that (a) uses shared\_initial\_candidate semantics exactly as registered, (b) archives raw model and validator traces for auditability, and (c) reports exact paired contingency counts together with the preregistered exact McNemar test and a paired bootstrap CI to aid interpretation in a small‑discordant sample.

\section{Methodology}
Overview and estimand.  We preregistered the study as cnsms2026-001-gvr-batfish-prereg-v1 and analysed the paired estimand labeled paired\_success\_rate\_difference\_guarded\_minus\_baseline under shared\_initial\_candidate semantics. Under those semantics the same single sampled initial model candidate (per task) is used to compute both outcomes: baseline = validator pass/fail on the initial candidate; guarded = final validator pass/fail after the allowed validator‑triggered repair (at most one LLM repair call). This design isolates the marginal effect of the permitted validator‑triggered repair for a fixed initial candidate rather than conflating effects of different initial samples or different first‑pass prompting strategies.

Task generation, inclusion/exclusion, and determinism.  The preregistered task bank deterministically produced 300 intent\ensuremath{\rightarrow}CLI pairs composed of public NetConfEval‑style examples and synthetic tasks generated by a seeded deterministic generator. Tasks are stratified into two vendor‑dialect clusters (150 Cisco‑like, 150 Juniper‑like). Generator IDs, seeds, and per‑task payloads (initial\_state, expected\_final\_state, required\_sequence, allowed\_touched\_fields, and workflow\_policies) are archived in the task manifest. Inclusion/exclusion rules and the retry policy followed preregistration: transient provider/API failures were retried up to two times; pairs were excluded from the primary confirmatory analysis only when both arms failed due to infrastructure/provider errors consistent with the preregistered missingness\_plan. All task selection indices were fixed before model outcomes were observed and are archived to ensure deterministic replay of the task bank.

Model, orchestration and recorded runtime parameters.  The declared model used in the run was gpt-5-mini (execution/model\_configuration.json). Archived execution metadata records max\_output\_tokens = 1024, maximum\_attempts\_per\_call = 1, provider = openai\_responses, and temperature = null (unset). Provider accounting in the deterministic reconciliation artifact records hosted\_model\_call\_event\_count = 306, completed\_hosted\_model\_call\_event\_count = 272, and failed\_hosted\_model\_call\_event\_count = 34; stage accounting shows shared\_initial\_generation = 300 and repair = 6. Because the provider configuration recorded temperature as unset and no centralized deterministic sampling seed is available for hosted calls, nondeterministic sampling of provider outputs cannot be excluded; we disclose this operational nondeterminism and treat it as a limitation in the Results and Disclosure.

Deterministic validator semantics and scoring rules.  The binary oracle was deterministic\_netops\_validator\_v5 and it applied a fixed battery of checks recorded in per‑episode validator traces. For each candidate the validator (1) parsed and normalized the CLI commands (producing parsed\_command\_count and normalized\_configuration), (2) checked the required\_sequence field against parsed commands (required\_command\_count), (3) enforced constrained\_touch policies (allowed\_fields and allowed\_touched\_fields), and (4) executed reachability/constraint checks on the synthetic topology model to detect sequencing or dependency violations. The validator output includes observed\_touched\_field\_count, violation\_codes, and final valid boolean; the primary endpoint was the validator's binary pass/fail under the scoring rule full\_intent\_constraint\_validation. Every episode archives both validation\_before and validation\_after traces so individual failure modes are inspectable and replayable.

Repair loop mechanics (bounded and archived).  The preregistered repair stage is strictly bounded to at most one automated LLM repair call per task when the initial candidate fails the validator. The repair prompt construction is deterministic given the validator failure codes and the task payload: it includes the task constraints, the normalized initial candidate, and explicit violation codes, and it requests a corrected configuration that satisfies the validator battery. Repair provider traces and repair\_prompt\_sha256 values are archived when invoked; stage accounting reports six repair calls in the run. Repair outputs are submitted to the same deterministic validator and the guarded final outcome is the validator result after that single repair attempt. Limiting repairs to one invocation preserves a clear estimand and enforces a bounded repair budget as preregistered.

Analysis procedures and exact inferential specification.  The prespecified primary inferential test is the exact McNemar two‑sided test on paired binary outcomes at \ensuremath{\alpha}=0.05. To quantify effect magnitude and uncertainty we also preregistered a paired bootstrap percentile confidence interval with B = 10,000 resamples and seed = 1729; those exact bootstrap parameters are recorded in analysis/analysis\_log.jsonl. Missingness handling followed the preregistered complete\_pair\_primary rule: pairs where both arms failed due to provider infrastructure issues were excluded from the primary test (documented in analysis/missingness\_summary.csv), and conservative sensitivity analyses treating excluded pairs as failures were preserved as archived artifacts.

Accounting, reconciliation and reproducibility artifacts.  Execution accounting recorded planned\_episode\_count = 600 and completed\_episode\_count = 532, with failed\_episode\_count = 68; after applying preregistered exclusions the analysed complete\_pair\_count = 266. Deterministic reconciliation artifacts verify internal consistency of paired counts and provider calls (provider and stage counts, cache behavior, and cross‑condition accounting) and assert that deterministic consistency checks passed. Key artifact categories archived for reproducibility include raw model responses, provider call traces, per‑episode scoring JSONs (validator traces), the task manifest, transformation manifests, model\_configuration.json, execution/raw\_results.jsonl, and the analysis result artifacts (condition\_summary.csv, paired\_contingency\_table.csv, missingness\_summary.csv, contamination\_summary.csv, and the paired\_difference\_figure.svg). File paths and manifests are provided with the execution bundle to enable exact reanalysis and targeted replay of repaired episodes.

Operational measurements attempted and absent.  The preregistration included secondary estimands for median\_repair\_iterations\_per\_task and median end‑to‑end latency; repair iterations were measurable and recorded (repair calls = 6; median repair iterations per task = 0). Per‑episode end‑to‑end latency (generation\ensuremath{\rightarrow}validation\ensuremath{\rightarrow}repair\ensuremath{\rightarrow}final validation) was not recorded in the archived per‑task artifacts and therefore the preregistered latency estimand could not be evaluated; this absence is noted as a preregistered deviation in the Limitations and drives a concrete reproducibility recommendation to include timing instrumentation in future runs.

Threats to validity embedded in the design.  The methodology preserves explicit distinctions required by the preregistration: internal validity is governed by shared\_initial\_candidate semantics (the estimand measures repair benefit for a fixed sample rather than differences from alternative first‑pass strategies), construct validity depends on validator rule coverage (the deterministic\_netops\_validator\_v5 battery defines which semantic violations are detectable), and external validity is limited by the two vendor clusters and synthetic topology scale. The preregistration's missingness rules and sensitivity analyses were applied exactly and all exclusions, retries, and provider failures are archived in the execution manifest to support transparent effect‑size interpretation.

\section{Results}
Execution accounting and sample construction. The preregistered run planned 600 model‑condition episodes (300 tasks \ensuremath{\times} 2 conditions). The orchestrated pipeline completed 532 episodes and recorded 68 failed episodes (execution\_manifest.json). After applying preregistered exclusion rules that remove pairs affected by both‑arm provider/infrastructure failures, the confirmatory analysis used 266 complete paired units (analysis/missingness\_summary.csv). Key provider‑call accounting from the deterministic reconciliation artifact: hosted\_model\_call\_event\_count = 306 total model call events, of which 272 completed successfully and 34 failed; cache was unused (cache\_hit\_count = 0, cache\_miss\_count = 272). Stage accounting shows that a shared initial generation was produced for every task (shared\_initial\_generation = 300) and the automated repair stage executed in 6 episodes (repair = 6). All execution‑level manifests and per‑episode traces are archived to permit replay and audit.

Primary paired outcomes and exact counts. The confirmatory paired contingency (guarded vs baseline) computed on the 266 complete pairs is shown in analysis/paired\_contingency\_table.csv and summarized here: n11 = 260 (both baseline and guarded pass), n10 = 5 (guarded pass, baseline fail), n01 = 0 (baseline pass, guarded fail), n00 = 1 (both fail). These margins yield observed success rates baseline = 260/266 = 0.9774436090 and guarded = 265/266 = 0.9962406015 with an absolute paired difference of 0.01879699248.

Inferential results: exact McNemar and paired bootstrap CI. Per the preregistered primary test, the exact McNemar two‑sided p‑value for the observed discordant counts (n10=5, n01=0) is p = 0.0625, which does not meet the preregistered \ensuremath{\alpha} = 0.05 decision rule for rejection. To quantify effect magnitude uncertainty we computed the preregistered paired bootstrap percentile 95\% CI with B = 10,000 resamples and seed = 1729 (analysis/analysis\_log.jsonl); the resulting CI is [0.0037593984962406013, 0.03759398496240601], narrowly excluding zero. Reporting both results preserves the preregistered decision rule while also conveying interval evidence: the discrete nature of McNemar's exact two‑sided test with few discordants can produce conservative p‑values even when bootstrap resampling indicates a small positive paired difference with limited sampling variability. Both inferential artifacts (exact p and bootstrap CI with documented seed and resamples) are archived.

Repair invocation behavior and representative repaired episodes. The repair stage executed for 6 of the 300 tasks (stage\_counts: repair = 6). Median repair iterations per task across the analysed sample is 0 (majority of tasks required no repair because the shared initial candidate already passed validation). Representative repaired episode: task-000008. The shared initial candidate for task-000008 (execution/responses/task-000008-shared-initial.txt) contained a truncated final token and failed validator checks (validation\_before showed parsed\_command\_count < required\_command\_count). The automated repair call produced a corrected candidate (execution/responses/task-000008-guarded.txt) whose validator trace (execution/scoring/task-000008-guarded.json) records validation\_after.valid = true with no violation codes. Per‑episode scoring JSONs record whether a repair was applied, the repair\_prompt reference, and the final normalized\_configuration accepted by the deterministic validator; all such per‑episode repair traces are archived for examination.

Contamination and sensitivity to missingness. The preregistered contamination detector flagged no exact public duplicates (analysis/contamination\_summary.csv empty), so no tasks were excluded from the primary analysis on contamination grounds. The preregistered missingness policy excluded 34 both‑arm failed episodes from the primary test; a sensitivity analysis that conservatively treats those 34 excluded pairs as failures (i.e., complete N = 300 with both arms set to failure where missing) yields baseline = 260/300 = 0.866667 and guarded = 265/300 = 0.883333, but retains the same discordant structure (n10=5, n01=0) leading to the same exact McNemar p = 0.0625. Thus, under the preregistered conservative missingness treatment the formal decision does not change; archived missingness\_summary.csv and deterministic\_reconciliation.json provide the complete accounting.

Supplementary quantitative artifacts and reproducibility. Analysis artifacts included in the evidence bundle and referenced above provide the reproducible numeric foundations for these results: analysis/condition\_summary.csv (baseline: 260 successes, 6 failures; guarded: 265 successes, 1 failure), analysis/paired\_contingency\_table.csv (n11, n10, n01, n00 table), analysis/missingness\_summary.csv (counts of excluded and failed episodes), and analysis/analysis\_log.jsonl (bootstrap settings: B=10,000, seed=1729). deterministic\_reconciliation.json documents provider call totals and stage counts that explain model\_calls\_used = 306 and repair\_calls = 6. All files and their hash manifests are archived to enable exact reanalysis or targeted follow‑up audits.

\section{Discussion}
The experiment produced a small absolute improvement in verifier pass‑rate when the allowed repair stage could operate on a shared initial candidate: guarded - baseline = 0.01879699248 (n11=260, n10=5, n01=0, n00=1). Practically, this pattern---few discordant pairs with all discordance in the guarded\ensuremath{\rightarrow}pass direction---indicates that (a) the initial LLM candidate was correct for most tasks in this corpus, and (b) the bounded repair stage occasionally corrected defects that the initial candidate incurred. The artifact accounting (repair calls = 6; hosted\_model\_call\_event\_count = 306) confirms repairs were rare relative to total generation activity, explaining the modest absolute gain.

From a statistical perspective the run highlights an important small‑sample inference nuance that directly affects operational interpretation. Exact McNemar (two‑sided) returned p = 0.0625 under the preregistered \ensuremath{\alpha}=0.05 decision rule, while a paired bootstrap percentile CI (B=10,000, seed=1729) produced a 95\% interval that narrowly excludes zero. These outcomes are consistent: with very few discordant pairs (here 5 vs 0), the exact two‑sided McNemar test is discrete and can be conservative; bootstrap resampling of paired differences summarizes empirical variability and can produce a CI excluding zero even when the exact discrete p slightly exceeds 0.05. For readers and practitioners this means that reporting both formal hypothesis outcomes and interval estimates (as archived here) gives a clearer picture: the effect magnitude is small but the CI suggests the effect is positive with narrow uncertainty in the observed corpus, whereas the exact hypothesis test under a strict \ensuremath{\alpha} threshold does not claim formal rejection.

Operationally, artifact evidence supports three practical inferences grounded in this run. First, when first‑pass generator accuracy is high (\ensuremath{\approx}97.7\% baseline here), bounded automated repair offers limited marginal benefit because few failures remain to correct; this is exactly what we observed (6 repair calls on 300 tasks). Second, the marginal utility of GVR will be larger when (i) first‑pass accuracy is lower, (ii) verifier coverage detects failure modes common in the target workload, or (iii) the repair mechanism is allowed more than one iteration---conditions not present in this bounded run. Third, the design of the verifier strongly conditions measured benefit: deterministic\_netops\_validator\_v5 enforces a finite, archived battery of syntactic/sequence/constraint/reachability checks; any semantic failure mode outside that battery is neither detected nor credited to the GVR pipeline. In short, GVR's measured improvement depends both on how often the generator errs and on whether the verifier can detect the errors that matter operationally.

Threats to inference and external validity remain artifact‑grounded. The preregistered shared\_initial\_candidate estimand isolates the effect of validator‑triggered repair for a fixed initial sample; it does not measure improvements obtainable by engineering the first pass (e.g., constrained prompting, constrained decoding, or multiple independent draws). Our execution manifest and reconciliation artifacts explicitly record stage\_counts (shared\_initial\_generation = 300, repair = 6) and the lack of independent first‑pass generations for separate arms, so extrapolating to comparisons that change initial sampling is unsupported by these artifacts. Missingness also reduced effective sample size: 34 both‑arm provider failures were excluded per preregistered rules, lowering precision and power; deterministic\_reconciliation.json and analysis/missingness\_summary.csv archive these counts. Finally, per‑episode latency was not recorded in per‑task artifacts and therefore the preregistered latency secondary estimand could not be evaluated---this absence constrains operational claims about responsiveness or real‑time applicability.

Design and research recommendations informed by these artifacts. (1) For stronger causal evidence about prompting or sampling strategies, future preregistrations should include additional independent arms that draw separate initial candidates (independent\_generation semantics) or include constrained‑first‑pass prompting; these require separate initial model calls per arm and differ from the shared\_initial\_candidate estimand used here. (2) To assess operational trade‑offs, pipeline runs must archive per‑episode timing (generation\ensuremath{\rightarrow}verification\ensuremath{\rightarrow}repair\ensuremath{\rightarrow}final validation) so latency and throughput can be measured; our provider accounting shows model call counts and failures but not per‑episode latencies. (3) Because verifier coverage materially shapes measured benefit, future studies should expand the verifier rule set and publish clear catalogs of detectable failure codes; the per‑episode validator traces archived in execution/scoring/*.json enable such audits and should be exploited to extend verifier coverage and to quantify verifier false negatives. (4) When baseline pass rates are high, consider targeting task corpora with intentionally lower baseline accuracy (e.g., out‑of‑distribution intents, noisier prompts, or reduced grounding) to increase discordant counts and thereby statistical power to detect repair benefits.

Concluding synthesis: the archived artifacts show that a bounded GVR pipeline can correct occasional deterministic‑validator failures for a fixed initial candidate, but when the generator's first‑pass accuracy is already high and the verifier's coverage is finite, the absolute operational improvement is small. The run's reconciled accounting, validator traces, repair provider traces, and contamination checks (no exact public duplicates flagged) provide a reproducible empirical substrate for these conclusions and for follow‑on experiments that expand sampling semantics, verifier coverage, or repair budgets.

\section{Limitations}
\begin{itemize}
\item Shared\_initial\_candidate semantics: both arms used the same single sampled initial candidate per preregistration; the estimand measures validator‑triggered repair benefit for that fixed initial sample and does not evaluate independent first‑pass generations or constrained first‑pass prompting strategies.
\item Missingness and sample reduction: planned N=300 pairs; after infrastructure/provider exclusions 266 pairs were complete and 34 both‑arm episodes were excluded per preregistered rules, reducing effective sample size and precision.
\item Small discordant counts: primary inference used exact McNemar with few discordants (n10=5, n01=0); conclusions are sensitive to inferential choice and limited power.
\item Validator coverage: deterministic\_netops\_validator\_v5 enforces a finite set of syntactic/semantic/reachability rules; semantic violations outside its coverage are possible and unmeasured.
\item Runtime nondeterminism and sampling: archived execution/model\_configuration.json records temperature = null and no deterministic seed; nondeterministic sampling cannot be excluded for recorded model calls.
\item H2 latency not evaluated: per‑episode end‑to‑end latency was not present in archived per‑task artifacts and therefore the preregistered latency bound (median <= 60 s) could not be assessed.
\item External validity: task corpus derives from public NetConfEval‑style examples and a seeded synthetic generator limited to two vendor‑dialect clusters; results do not imply universal benefit across other vendors, larger topologies, or production deployments.
\end{itemize}

\section{Conclusion}
We executed a preregistered paired experiment that evaluated a GVR pipeline (gpt‑5‑mini + deterministic\_netops\_validator\_v5 + <=1 automated repair) on intent\ensuremath{\rightarrow}CLI tasks under shared\_initial\_candidate semantics. Across 266 analysed pairs the guarded pipeline improved validator pass rate by 0.01879699248 but did not meet the preregistered exact‑McNemar \ensuremath{\alpha}=0.05 decision rule (p=0.0625); a paired bootstrap CI (B=10,000, seed=1729) narrowly excludes zero. Repair calls were rare (6/300) and median repair iterations per task = 0. Per‑episode latency was not archived and could not be assessed. All execution and analysis artifacts (raw responses, validator traces, provider call accounting, and reconciliation manifests) are archived in the evidence bundle to support reanalysis and follow‑on work.

References
[1] C. Wang, M. Scazzariello, A. Farshin, S. Ferlin, D. Kostić, and M. Chiesa, "NetConfEval: Can LLMs Facilitate Network Configuration?," in Proc. ACM, 2024, doi: 10.1145/3656296.
[2] S. Y. Long, J. Tan, B. Mao, F. Tang, Y. Li, M. Zhao, and N. Kato, "A Survey on Intelligent Network Operations and Performance Optimization Based on Large Language Models," IEEE Commun. Surveys \& Tutorials, 2025, doi: 10.1109/comst.2025.3526606.
[3] R. Cadenas, "Convergent Correctness in Stochastic Code Generation: A Generate‑Verify‑Repair Architecture for Deterministic Validation of Autonomous AI Coding Agents in Regulated Environments," SSRN, 2026, doi: 10.2139/ssrn.6754899.
[4] M. Brown, A. Fogel, D. Halperin, V. Heorhiadi, R. Mahajan, and T. Millstein, "Lessons from the evolution of the Batfish configuration analysis tool," Proc., 2023, doi: 10.1145/3603269.3604866.
[5] R. Mondal, A. Tang, R. Beckett, T. Millstein, and G. Varghese, "What do LLMs need to Synthesize Correct Router Configurations?," Proc., 2023, doi: 10.1145/3626111.3628194.

\section*{Disclosure Statement}
Disclosure (concise): Declared model: gpt-5-mini (execution/model\_configuration.json archived; declared\_model\_version = "gpt-5-mini"). Orchestration adapter: hosted\_netops\_gvr\_v1. Deterministic validator: deterministic\_netops\_validator\_v5. Preregistration: cnsms2026-001-gvr-batfish-prereg-v1 (manifest SHA‑256: 46a22349\allowbreak{}2f760950\allowbreak{}b3d06f47\allowbreak{}5129ce21\allowbreak{}e52e296c\allowbreak{}4c2a1e09\allowbreak{}df3ccd50\allowbreak{}bec9f891). Initial master prompt immutable reference: provenance/master\_prompt.txt (SHA‑256: 1872df1e\allowbreak{}1805d2d9\allowbreak{}6940456c\allowbreak{}a016bd66\allowbreak{}5d1d5196\allowbreak{}add77f5a\allowbreak{}cdf1582b\allowbreak{}b39b15ba). Human role and autonomy boundary: before the autonomy lock a human Research Director selected the topic family and approved pre‑lock infrastructure development; after the autonomy lock the registered pipeline autonomously performed literature selection, preregistration, experimental execution, analysis, manuscript drafting, AI peer review simulation, and revision. Nondeterminism disclosure: archived execution/model\_configuration.json records temperature = null (unset) and no deterministic sampling seed is archived; therefore nondeterministic sampling of model outputs cannot be excluded for the recorded provider calls. Execution and analysis artifacts, logs, and hash manifests are archived in the evidence bundle referenced in the preregistration.

\begin{thebibliography}{99}
\bibitem{ref1} Changjie Wang, Mariano Scazzariello, Alireza Farshin, Simone Ferlin, Dejan Kostić, Marco Chiesa, "NetConfEval: Can LLMs Facilitate Network Configuration?", 2024, 10.1145/3656296
\bibitem{ref2} S.Y. Long, Jingjing Tan, Bomin Mao, Fengxiao Tang, Yangfan Li, Ming Zhao, Nei Kato, "A Survey on Intelligent Network Operations and Performance Optimization Based on Large Language Models", 2025, 10.1109/comst.2025.3526606
\bibitem{ref3} Rafael Cadenas, "Convergent Correctness in Stochastic Code Generation: A Generate-Verify-Repair Architecture for Deterministic Validation of Autonomous AI Coding Agents in Regulated Environments", 2026, 10.2139/ssrn.6754899
\bibitem{ref4} Matt Brown, Ari Fogel, Daniel Halperin, Victor Heorhiadi, Ratul Mahajan, Todd Millstein, "Lessons from the evolution of the Batfish configuration analysis tool", 2023, 10.1145/3603269.3604866
\bibitem{ref5} Rajdeep Mondal, Alan Tang, Ryan Beckett, Todd Millstein, George Varghese, "What do LLMs need to Synthesize Correct Router Configurations?", 2023, 10.1145/3626111.3628194
\end{thebibliography}

\end{document}
